\section{Extending the household stage}\label{sec:roadmap}

Section~\ref{sec:results-secondstage} runs the energy episode through
PolicyEngine UK on the enhanced Family Resources Survey. Extending the
same stage to the remaining four episodes is a parameter reform per
episode rather than new machinery, and four things change when it is
done.

First, the consumption channel moves from decile tables to household
records: the enhanced FRS carries imputed household consumption across
the twelve COICOP divisions plus the sub-items that matter most for
trade shocks (petrol and diesel separately, in pounds and litres;
electricity and gas separately; bus fares), imputed by quantile
regression forest from the LCFS 2023--24 and calibrated to
administrative energy and fuel totals. Each episode's price vector
therefore prices each household's own basket, jointly with that
household's benefit entitlements, the covariance between exposure and
cushioning that decile tables cannot see. Tariff-line and
category-level first stages must be pre-aggregated to those roughly
sixteen categories, and the imputation freezes budget-share
\emph{patterns} at the 2023--24 donor vintage across simulated years, a
limitation the stage will state. Second, the earnings channels run through the companion study's
displacement machinery on the same households, so the two channels of a
given episode land on a common population and their incidence can be
compared household-by-household rather than distribution-by-
distribution. Third, the rulebook axis becomes exact: PolicyEngine's parameter
history carries dated, legislatively cited values through the whole
window (the Covid-era Universal Credit uplift and its October 2021
removal, the March 2022 fuel-duty cut, the 2022 and 2023 upratings),
its own test suite exercises the 2020--2023 rule years, and
counterfactual-rulebook runs, an earlier year's parameters applied to
a later survey vintage, are native to the architecture. Each episode
is therefore scored under the rules in force at its date, under the
preceding year's rules (isolating uprating lag), and with the
discretionary interventions switched off (isolating the automatic
component), the three-way decomposition of
Section~\ref{sec:framework} computed rather than approximated. Two
verified limits shape the earlier years: no 2021-vintage enhanced-FRS
data year exists (the earliest survey base is 2022--23), so 2021
rulebooks are applied to later households with incomes re-deflated;
and the survey's Universal Credit/legacy-benefit split follows
reported receipt rather than the historical rollout, so 2021--22
rule-year results understate legacy-benefit interactions and are
restricted to the tax and UC-parameter channels. Fourth, the poverty and inequality outputs gain their
correct household equivalisation and the Exchequer effects their full
statutory coverage.

The build order follows the data dependencies. The consumption-channel
implementation extends PolicyEngine's existing consumption variables
(used today for carbon and VAT analysis) with the episode price
vectors, a parameter reform per episode, no new imputation. The
earnings channels are re-runs of existing machinery with episode
calibrations. The rulebook runs require only parameter-year switches
plus explicit flags for the discretionary layers (the Energy Price
Guarantee being the largest). The licensed inputs are the same FRS/EFRS
files the companion study uses; nothing else is gated.

Two methodological upgrades are available at the same time and are
specified here. First, the fixed-weight decile indices of the first
pass can be replaced by non-homothetic distributional cost-of-living
indices, which \citet{hochmuth2025} construct from exactly the inputs
this project already has, aggregate prices, expenditure shares and a
single cross-sectional expenditure distribution, answering the
standard objection that fixed baskets misprice large relative-price
movements. Second, and more consequentially, the household-record
stage is the level at which the within-decile dispersion documented by
\citet{borusyakjaravel2024} and \citet{levell2025} becomes visible: the
same machinery that computes a decile gradient computes the full
distribution of exposure within it, which is where most of the welfare
variance lives.

Two honest limits will survive the full stage. The consumption
imputation carries its own error, inherited from the LCFS donor
models, and imputation error correlated with benefit status would bias
the exposure-cushioning covariance; the stage will report the
imputation diagnostics alongside the results. And the rulebook
decomposition attributes to ``discretion'' whatever the parameter
history records as legislated change, which bundles genuine shock
response with policy that would have happened anyway, a labelling
choice, stated as such, not an identification claim.
